\documentclass[12pt]{letter}
\usepackage[margin=2.5cm]{geometry}
\usepackage{hyperref}
\usepackage{amsmath,amssymb}

\signature{Jian Zhou\\[2pt]
\textit{Independent Researcher}\\[2pt]
\texttt{jackzhou.sci@gmail.com}}

\address{Jian Zhou\\Independent Researcher\\jackzhou.sci@gmail.com}

\date{\today}

\begin{document}

\begin{letter}{Editorial Office\\
\textit{Journal of Mathematical Physics}\\
AIP Publishing}

\opening{Dear Editor,}

I am pleased to submit the manuscript entitled
\textbf{``Stokes Multipliers of $x^{2M}$ Anharmonic Oscillators:
Exact Results for $M = 2$--$11$''}
for consideration as an Article in the \textit{Journal of Mathematical Physics}.

This paper presents a systematic high-precision computation of the
Stokes multipliers~$C_M$ for the one-parameter family of anharmonic
oscillators $H = p^2/2 + x^2/2 + g\,x^{2M}$, covering $M = 2$
through~$11$. The main results are:

\begin{itemize}
\item \textbf{Three new exact closed forms} for the Stokes multipliers
  $C_3$, $C_5$, and~$C_7$, expressed in terms of Gamma values at
  rational arguments and~$\pi$. To our knowledge, these have not
  appeared previously in the literature.

\item \textbf{Strong numerical evidence} (30-digit PSLQ exclusion) that
  $C_4$ admits no such closed form, identifying the octic oscillator as
  a unique anomaly.

\item A \textbf{number-theoretic pattern} connecting closed-form
  existence to the Euler totient function~$\varphi(M{-}1)$ and the
  algebraic independence structure of Gamma values
  (Rohrlich--Lang conjecture).
\end{itemize}

The exact values of $C_5$ and $C_7$ provide new non-trivial benchmarks
for resurgent analyses of higher anharmonic oscillators, a topic of
active interest in mathematical physics following the work of
Zinn-Justin--Jentschura, Dunne--\"Unsal, and others. The unexpected
connection between Stokes multipliers and the arithmetic of Gamma
functions may be of independent interest to researchers working on
periods and transcendental number theory.

All computational code and data are publicly available at
\url{https://github.com/JackZH26/stokes-multipliers}, and the
manuscript is posted as a Figshare preprint
(DOI: \href{https://doi.org/10.6084/m9.figshare.31796332.v2}{10.6084/m9.figshare.31796332.v2}).

This manuscript has not been submitted elsewhere and all work is
original. I confirm that I have no conflicts of interest to declare.

I believe this work falls within the scope of JMP's coverage of
quantum mechanics, spectral theory, and mathematical methods in
physics. I would suggest the following as potential referees with
expertise in this area:

\begin{itemize}
\item Prof.\ Gerald V.\ Dunne (University of Connecticut) ---
  resurgence and non-perturbative methods in quantum mechanics.
\item Prof.\ Ulrich D.\ Jentschura (Missouri University of Science
  and Technology) --- multi-instanton expansions, anharmonic oscillators.
\item Prof.\ Ovidiu Costin (Ohio State University) --- exact WKB,
  Borel summability, Stokes phenomena.
\end{itemize}

Thank you for considering this submission.

\closing{Sincerely,}

\end{letter}
\end{document}
